\documentclass[letterpaper,11pt]{article}
\usepackage[T1]{fontenc}
\usepackage[utf8]{inputenc}
\usepackage[hidelinks]{hyperref}
\usepackage{xcolor}
\usepackage[margin=1in]{geometry}
\usepackage[default]{sourcesanspro}
\definecolor{ResumeNavy}{HTML}{18324A}
\definecolor{ResumeSlate}{HTML}{5B6573}
\color{black}
\setlength{\parskip}{8pt}
\setlength{\parindent}{0pt}
\hypersetup{colorlinks=true, urlcolor=ResumeNavy}
\begin{document}
\begin{center}
  {\LARGE \bfseries \textcolor{ResumeNavy}{Deva Anand}} \\[3pt]
  {\small
  \href{mailto:devaanand@umass.edu}{devaanand@umass.edu} \enspace\textbar\enspace
  (413) 315-1715 \enspace\textbar\enspace
  \href{https://github.com/Deva-1903}{github.com/Deva-1903} \enspace\textbar\enspace
  \href{https://linkedin.com/in/devaaa/}{linkedin.com/in/devaaa}}
\end{center}
\vspace{6pt}

June 8, 2026

\vspace{2pt}
PathAI \\
Machine Learning Intern/Co-op Hiring Team \\
Boston, MA

\vspace{6pt}
Dear Hiring Team,

I'm applying for the Machine Learning Intern/Co-op role. I'm a Computer Science master's student at UMass Amherst, and PathAI stands out to me because the work sits exactly where I want to be: applying machine learning to medical images where the output is a diagnosis, not a metric on a leaderboard.

That's not abstract for me. My first publication (IEEE CONIT 2023, first author) applied deep transfer learning to neuroimaging for Alzheimer's classification, a computer-vision approach to disease diagnosis. In my graduate coursework I went deeper on the fundamentals: I trained and compared CNN and ResNet architectures for image classification on CIFAR-10, implemented a diffusion model and a normalizing flow from scratch in JAX, and even built a reverse-mode autodiff engine by hand and validated it against JAX to machine precision. I like understanding models from the ground up, which I think matters on a team that has to trust what its models are doing.

The other half of your work is rigorous evaluation, and that's where I've spent a lot of energy. I built an evaluation-improvement agent that only accepts changes that increase a benchmark's discriminative power, with guardrails against gaming the holdout. And in a recent interpretability study I designed proper controls and reported a clean null result, with bootstrap confidence intervals and a McNemar's test, rather than the finding I'd hoped for. Careful experimental design and honest analysis are the parts of ML I take most seriously, which is why your emphasis on rigorous science resonated.

I also care about the patient-outcomes mission specifically. Before grad school I spent two years as an engineer at Wysa, a clinical mental-health company, where I built the ML training-data pipeline the AI team used to label clinical data for production models, and owned reliability on a system serving real NHS patients. Working where mistakes actually affect people is something I want more of, not less.

I know my recent depth is broader ML rather than digital pathology specifically, but I'm a fast learner, I have the medical-imaging and evaluation foundation, and I'd be genuinely excited to contribute. Thank you for considering my application.

\vspace{8pt}
Sincerely, \\
Deva Anand
\end{document}
